\begin{itemize}[leftmargin=*]
  \setlength{\itemsep}{0.35em}
  \item \textbf{Custo médio/máquina (comparativo justo):} LIEBHERR custa \textbf{42.8\%} a mais (R\$ 124,072/máq vs R\$ 86,871/máq; parque: 7 vs 16 máquinas).
  \item \textbf{Volume absoluto (contexto):} HYUNDAI tem maior gasto total (R\$ 1,399,720 vs R\$ 870,203) — reflexo do tamanho do parque.
  \item \textbf{2025 — HYUNDAI:} volume R\$ 1,220,373 (custo médio R\$ 75,649/máquina).
  \item \textbf{2025 — LIEBHERR:} volume R\$ 733,840 (custo médio R\$ 104,834/máquina).
  \item \textbf{2026 — HYUNDAI:} volume R\$ 179,347 (custo médio R\$ 11,209/máquina).
  \item \textbf{2026 — LIEBHERR:} volume R\$ 136,362 (custo médio R\$ 19,237/máquina).
  \item \textbf{Tipo dominante (R\$/máq):} LIEBHERR $\rightarrow$ 7.1.2 MANUTENÇÃO DE VEÍCULOS/MAQUINAS (R\$ 52,641/máq); HYUNDAI $\rightarrow$ 7.1.2 MANUTENÇÃO DE VEÍCULOS/MAQUINAS (R\$ 46,949/máq).
  \item \textbf{Mix manutenção (7.1.2):} LIEBHERR 42.4\% vs HYUNDAI 54.1\% do volume.
  \item \textbf{Mix combustível (7.1.22+7.1.4):} LIEBHERR 25.2\% vs HYUNDAI 35.6\%.
  \item \textbf{Mix peças (7.1.1):} LIEBHERR 30.5\% vs HYUNDAI 8.2\%.
  \item \textbf{Concentração LIEBHERR:} top 3 máquinas = 82\% do gasto (líder: EHL0043).
  \item \textbf{Concentração HYUNDAI:} top 3 máquinas = 38\% do gasto (líder: EHH0044).
  \item \textbf{Unidade com maior custo:} Maringá — R\$ 556,716 (5 máq.; R\$ 111,343/máq).
  \item \textbf{Top 3 unidades:} Maringá (R\$ 556,716); Dourados (R\$ 429,208); Prudente (R\$ 349,489) — 59\% do gasto do parque.
  \item \textbf{Maior intensidade (R\$/máq):} Dourados — R\$ 143,069/máq (3 máq.).
  \item \textbf{Na unidade líder (Maringá):} LIEBHERR R\$ 348,601 vs HYUNDAI R\$ 208,115 (predomina LIEBHERR).
\end{itemize}
